% ============================================================
%  吴宛幸  个人简历 (中文版)
% ============================================================
\documentclass[10pt, a4paper]{article}

% ── Encoding & CJK Support ──────────────────────────────────
\usepackage[UTF8, heading=false]{ctex}
\usepackage[T1]{fontenc}
\usepackage{palatino}
\usepackage{mathpazo}
\usepackage{microtype}

% ── Page Layout ──────────────────────────────────────────────
\usepackage[top=1.2cm, bottom=1.2cm, left=1.5cm, right=1.5cm]{geometry}
\setlength{\parindent}{0pt}
\setlength{\parskip}{0pt}
\linespread{1.0}\selectfont

% ── Packages ─────────────────────────────────────────────────
\usepackage{xcolor}
\usepackage{hyperref}
\usepackage{enumitem}
\usepackage{titlesec}
\usepackage{tabularx}
\usepackage{array}
\usepackage{fancyhdr}
\usepackage{tikz}
\usepackage{etoolbox}
\usepackage{graphicx}

% ── Colors ───────────────────────────────────────────────────
\definecolor{headcol}{RGB}{30, 30, 30}
\definecolor{rulecol}{RGB}{180, 50, 30}
\definecolor{secline}{RGB}{200, 200, 200}
\definecolor{linkcol}{RGB}{40, 80, 160}
\definecolor{dategray}{RGB}{100, 100, 100}

% Tag colors
\definecolor{tagroute}{RGB}{130, 50, 160}
\definecolor{taggenomics}{RGB}{40, 140, 90}
\definecolor{tagrl}{RGB}{180, 100, 20}

\hypersetup{
  colorlinks = true,
  urlcolor   = linkcol,
  linkcolor  = linkcol,
  pdfauthor  = {吴宛幸},
  pdftitle   = {吴宛幸，个人简历},
}

% ── Section Headers ──────────────────────────────────────────
\titleformat{\section}
  {\large\bfseries\color{headcol}}
  {}{0em}{}
  [\vspace{2pt}{\color{secline}\rule{\linewidth}{0.4pt}}\vspace{-2pt}]
\titlespacing*{\section}{0pt}{8pt}{3pt}

% ── Badge Tags ───────────────────────────────────────────────
\newcommand{\badge}[2]{%
  \tikz[baseline=(x.base)]{
    \node[inner xsep=3.5pt, inner ysep=1.6pt, rounded corners=2.5pt,
          draw=#1, line width=0.45pt,
          text=#1, font=\fontsize{6}{7}\selectfont\sffamily]
      (x) {#2};
  }%
}

\newcommand{\troute}  {\badge{tagroute}{routing}}
\newcommand{\tgenomics}{\badge{taggenomics}{genomics}}
\newcommand{\trl}     {\badge{tagrl}{RL}}

% ── Bullet Lists ─────────────────────────────────────────────
\setlist[itemize,1]{
  leftmargin = 1.3em,
  label      = {\small\textbullet},
  itemsep    = 0pt,
  topsep     = 2pt,
  parsep     = 0pt,
}

% ── Helpers ──────────────────────────────────────────────────
\newcommand{\unilogo}[1]{%
  \raisebox{-0.5pt}{\includegraphics[height=8pt]{assets/logos/cv/#1}}%
}
\newcommand{\me}[1]{\textbf{#1}}
\newcommand{\eqstar}{$^{*}$}
\newcommand{\corr}{$^{\dagger}$}

\def\pubskip{2pt}

% Single publication block
\newcommand{\pub}[4]{%
  \vspace{\pubskip}%
  \noindent\textbf{#1}\ifstrempty{#2}{}{\enskip #2}%
  \begin{itemize}[topsep=0pt]
    \item #3
    \item \textit{#4}
  \end{itemize}%
}

% Education / Experience entry
\newcommand{\entry}[2]{%
  \noindent
  \begin{tabularx}{\linewidth}{@{}X r@{}}
    \textbf{#1} & {\small\color{dategray}#2}
  \end{tabularx}%
  \par\vspace{0pt}%
}

% ── Header / Footer ──────────────────────────────────────────
\pagestyle{fancy}
\fancyhf{}
\renewcommand{\headrulewidth}{0pt}
\fancyfoot[R]{\footnotesize\color{dategray}\thepage}

% ============================================================
\begin{document}
\thispagestyle{empty}

% ── NAME BLOCK ───────────────────────────────────────────────
\noindent
\begin{tabularx}{\linewidth}{@{}X r@{}}
  \raisebox{-4pt}{%
    {\fontsize{28}{32}\selectfont\bfseries 吴宛幸}%
    \quad{\large\color{dategray}2027届}%
  }
  &
  \begin{tabular}[b]{@{}r@{}}
    \small\href{mailto:wuwanxing3574@gmail.com}{wuwanxing3574@gmail.com}\\[2pt]
    \small\href{https://github.com/wwx777}{GitHub}
    \quad \href{https://wwx777.github.io}{个人主页}
  \end{tabular}
\end{tabularx}

\vspace{4pt}
{\color{rulecol}\rule{\linewidth}{1.2pt}}
\vspace{1pt}

% ── ABOUT ────────────────────────────────────────────────────
\section{个人简介}

\textbf{巴黎综合理工学院（IP Paris）}计算机科学硕士在读，本科毕业于\textbf{南方科技大学}（数据科学与大数据技术），现为 \textbf{SUSTech-NLP} 科研助理（导师：陈冠华教授）。

\vspace{2pt}
我的研究方向是\textbf{异构能力的 LLM 应该如何协作}。我从 per-query 路由切入：在 \textbf{RouterXBench} 中，我设计了一套将 router 内在能力与场景适配性解耦的三轴评估框架。同样的协作问题在 agent 部署到新环境、需要在多步任务中适应时再次出现；在 \textbf{Scope-then-Cover} 中，我将其实例化为单一 explore action：强 actor 提名小候选子空间，冻结的廉价 worker 并行验证；该方案在 test-time adaptation (TTA) 与 online experiential learning (OEL) 两种场景下均有效。

\vspace{2pt}
一作 / 共同一作论文 2 篇，均在审。

% ── EDUCATION ────────────────────────────────────────────────
\section{教育经历}

\entry{\unilogo{ipp.png}\enskip 巴黎综合理工学院（Institut Polytechnique de Paris）}{法国巴黎}
\textit{计算机科学硕士} \hfill \textit{2025.09 -- 2027.09}
\begin{itemize}
  \item \textbf{ReflexSQL}：基于探索与记忆的 Text-to-SQL agent，验证于 Spider~2.0-Lite \quad [\href{https://wwx777.github.io/assets/reflexsql-poster.pdf}{海报}]。
\end{itemize}

\vspace{2pt}
\entry{\unilogo{sustech.png}\enskip 南方科技大学}{深圳}
\textit{数据科学与大数据技术，理学学士} \hfill \textit{2021.09 -- 2025.07}
\begin{itemize}
  \item 本科毕业论文：大小模型协同路由（在 \textbf{RouterXBench} 中延续）。
\end{itemize}

% ── PUBLICATIONS ─────────────────────────────────────────────
\section{学术论文}

\pub{Towards Fair and Comprehensive Evaluation of Routers in Collaborative LLM Systems}{}
    {\me{Wanxing Wu}, He Zhu, Yixia Li, \ldots, Guanhua Chen\corr}
    {在审，2026 \quad[\href{https://arxiv.org/abs/2602.11877}{arXiv:2602.11877}]}

\pub{Can Weak Models Explore? Scope Then Cover for LLM Agent Adaptation}{}
    {\me{Wanxing Wu}\eqstar, He Zhu\eqstar, Guanhua Chen\corr
     \quad ($^*$同等贡献)}
    {在审，2026}

% ── RESEARCH EXPERIENCE ──────────────────────────────────────
\section{科研经历}

\entry{\unilogo{sustech.png}\enskip SUSTech-NLP 实验室}{深圳}
\textit{科研助理}，导师：陈冠华教授 \hfill \textit{2025.01 -- 至今}

异构能力 LLM 协作：per-query 路由与 agent 长任务中的 per-step 分工。

\vspace{2pt}
\entry{华大基因（BGI）智能医学研究院（IIMR）}{深圳}
\textit{机器学习实习生（AI for Genomics）} \hfill \textit{2024.07 -- 2025.01}

cfDNA 特征表示 $+$ 集成学习用于结直肠癌早筛，灵敏度 \textbf{94.88\%}。

% ── HONORS ───────────────────────────────────────────────────
\section{荣誉与奖项}

\begin{itemize}[itemsep=0pt, topsep=1pt]
  \item \textbf{优秀毕业生}，南方科技大学致仁书院（2025.07）
  \item \textbf{优秀毕业论文}，Top 5\%，南方科技大学（2025.07）
  \item 优秀学生奖学金，南方科技大学（2024.10）
\end{itemize}

% ── SKILLS ───────────────────────────────────────────────────
\section{技能}

\noindent\textbf{大模型与机器学习：} 推理框架（vLLM、SGLang），Agentic-RL 工作流，SFT/RL 训练。\\[3pt]
\noindent\textbf{编程与工具：} Python、PyTorch、Docker、Linux、Git、Shell 脚本。\\[3pt]
\noindent\textbf{语言：} 英语（IELTS 6.5，CET-4，CET-6）；中文（母语）。

\end{document}
